% \begin{figure}[htbp]
% \centering
% \begin{tikzpicture}[
%     x=1cm, y=1cm,
%     font=\small,
%     bar/.style={draw, rectangle, rounded corners=2pt, align=center, inner sep=2pt, font=\scriptsize},
%     callout/.style={draw=black!25, fill=black!2, rectangle, rounded corners=2pt, align=center, font=\scriptsize, inner sep=4pt}
% ]

% % Total width budget: ~7.2cm (Fits standard 1-column layout perfectly)

% % =========================================================
% % PART (a): Intra-Audio Chunk Graph Splits
% % =========================================================
% \node[anchor=west, font=\bfseries\small] at (0, 7.8) {(a) Graph Chunks Inside One Audio Chunk};

% % 30s Audio Input Context
% \node[bar, fill=inputbg, draw=inputborder, minimum width=7.2cm, minimum height=0.45cm] at (3.6, 7.2) {
%     \textbf{Audio Input $c$} \textit{(15s -- 45s)}
% };

% % 20s Retained Center
% \node[bar, fill=retainedbg, draw=retainedborder, minimum width=4.8cm, minimum height=0.45cm] at (3.6, 6.55) {
%     \textbf{Retained Output Center $c$} \textit{(20s -- 40s)}
% };

% % 2x 10s Graph Chunks
% \node[bar, fill=graphbg1, draw=graphborder1, minimum width=2.3cm, minimum height=0.45cm] at (2.35, 5.9) {
%     [20--30s] \textbf{($\pi_c$)}
% };
% \node[bar, fill=graphbg1, draw=graphborder1, minimum width=2.3cm, minimum height=0.45cm] at (4.85, 5.9) {
%     [30--40s] \textbf{($\pi_c$)}
% };

% \node[font=\scriptsize, color=black!70] at (3.6, 5.35) {
%     \textit{All graph sub-chunks share the identical speaker map $\pi_c$.}
% };

% % Separator line
% \draw[black!20, thick] (0, 5.0) -- (7.2, 5.0);

% % =========================================================
% % PART (b): Inter-Audio Chunk Graph Splits
% % =========================================================
% \node[anchor=west, font=\bfseries\small] at (0, 4.6) {(b) Graph Chunks Across Audio Chunks};

% % Overlapping Audio Context Inputs
% \node[bar, fill=inputbg!70, draw=inputborder!70, minimum width=4.2cm, minimum height=0.4cm] at (2.1, 4.0) {
%     Audio Input $c$ \textit{[-5s -- 25s]}
% };
% \node[bar, fill=inputbg!70, draw=inputborder!70, minimum width=4.2cm, minimum height=0.4cm] at (5.1, 3.45) {
%     Audio Input $c+1$ \textit{[15s -- 45s]}
% };

% % Non-overlapping Contiguous Retained Centers
% \node[bar, fill=retainedbg, draw=retainedborder, minimum width=3.4cm, minimum height=0.45cm] at (1.7, 2.8) {
%     Retained $c$ \textit{[0s -- 20s]}
% };
% \node[bar, fill=retainedbg, draw=retainedborder, minimum width=3.4cm, minimum height=0.45cm] at (5.3, 2.8) {
%     Retained $c+1$ \textit{[20s -- 40s]}
% };

% % Graph Chunks laid out directly under their respective centers
% \node[bar, fill=graphbg1, draw=graphborder1, minimum width=1.6cm, minimum height=0.45cm] at (0.8, 2.15) {
%     [0--10s] \textbf{$\pi_c$}
% };
% \node[bar, fill=graphbg1, draw=graphborder1, minimum width=1.6cm, minimum height=0.45cm] at (2.6, 2.15) {
%     [10--20s] \textbf{$\pi_c$}
% };

% \node[bar, fill=graphbg2, draw=graphborder2, minimum width=1.6cm, minimum height=0.45cm] at (4.4, 2.15) {
%     [20--30s] \textbf{$\pi_{c+1}$}
% };
% \node[bar, fill=graphbg2, draw=graphborder2, minimum width=1.6cm, minimum height=0.45cm] at (6.2, 2.15) {
%     [30--40s] \textbf{$\pi_{c+1}$}
% };

% % Mapping Transition Box
% \node[callout, minimum width=7.2cm] at (3.6, 1.0) {
%     \textbf{Speaker Mapping Transition Across Audio Boundary:}\\
%     \textbf{Chunk $c$ ($\pi_c$):} Alice $\to$ Slot 0, Bob $\to$ Slot 1\\
%     \textbf{Chunk $c+1$ ($\pi_{c+1}$):} Alice $\to$ \textbf{Slot 1}, Bob $\to$ \textbf{Slot 0}\\
%     \vspace{1pt}
%     \textit{Retained graph chunks step cleanly from $\pi_c$ to $\pi_{c+1}$ at 20s.}
% };

% \end{tikzpicture}
% \caption{Structural breakdown showing how Viterbi graph chunks inherit speaker maps within a single audio chunk (a) versus across contiguous audio chunk boundaries (b).}
% \label{fig:speaker_tracking_parallel}
% \end{figure}


\begin{figure}[htbp]
\centering
\begin{tikzpicture}[
    x=1cm, y=0.65cm,
    font=\small,
    bar/.style={draw, rectangle, rounded corners=2pt, align=center, inner sep=2pt, font=\scriptsize}
]

% Total width: 7.0cm (Fits 1-column IEEE/ACL layout)

% 1. Contiguous Retained Audio Output Centers
\node[bar, fill=retainedbg, draw=retainedborder, minimum width=3.4cm, minimum height=0.42cm] at (1.7, 2.5) {
    \textbf{Retained Audio $c$} \textit{[0--20s]}
};
\node[bar, fill=retainedbg, draw=retainedborder, minimum width=3.4cm, minimum height=0.42cm] at (5.3, 2.5) {
    \textbf{Retained Audio $c+1$} \textit{[20--40s]}
};

% 2. Graph Sub-Chunks
% Inside Chunk c (Intra-chunk: Both use \pi_c)
\node[bar, fill=graphbg1, draw=graphborder1, minimum width=1.6cm, minimum height=0.42cm] at (0.8, 1.6) {
    0--10s \textbf{($\pi_c$)}
};
\node[bar, fill=graphbg1, draw=graphborder1, minimum width=1.6cm, minimum height=0.42cm] at (2.6, 1.6) {
    10--20s \textbf{($\pi_c$)}
};

% Inside Chunk c+1 (Inter-chunk boundary at 20s -> switches to \pi_{c+1})
\node[bar, fill=graphbg2, draw=graphborder2, minimum width=1.6cm, minimum height=0.42cm] at (4.4, 1.6) {
    20--30s \textbf{($\pi_{c+1}$)}
};
\node[bar, fill=graphbg2, draw=graphborder2, minimum width=1.6cm, minimum height=0.42cm] at (6.2, 1.6) {
    30--40s \textbf{($\pi_{c+1}$)}
};

% 3. Minimal Annotations
\draw[thick, graphborder1, <->] (0.0, 0.95) -- (3.4, 0.95) 
    node[midway, below, font=\tiny\bfseries, color=graphborder1!90!black] {Intra-Chunk: Shared $\pi_c$};

\draw[thick, graphborder2, <->] (3.6, 0.95) -- (7.0, 0.95) 
    node[midway, below, font=\tiny\bfseries, color=graphborder2!90!black] {Inter-Chunk: Transition to $\pi_{c+1}$};

% 4. One-Line Permutation Legend
\node[draw=black!20, fill=black!2, rectangle, rounded corners=2pt, font=\tiny, inner sep=2.5pt] at (3.5, -0.1) {
    \textbf{Speaker Slots:} $\pi_c = \{\text{Alice}\to0, \text{Bob}\to1\} \;\;\longrightarrow\;\; \pi_{c+1} = \{\text{Alice}\to1, \text{Bob}\to0\}$
};

\end{tikzpicture}
\caption{Unified speaker mapping mechanism: sub-chunks within audio chunk $c$ inherit $\pi_c$, while crossing the 20s boundary transitions directly to $\pi_{c+1}$.}
\label{fig:speaker_tracking_ultra_minimal}
\end{figure}